\section{Two precise rigidity statements}
\label{clc:sec:rigidity}

\subsection{When can the independent orders add purely?}
\begin{theorem}[Affine-center alignment is necessary and sufficient]
\label{clc:thm:alignment}
Fix positive scales and centers with a common strip. The convolution
function $\mu\mapsto\UU_{\bq,\bA}(\bs;\mu)$ depends on the orders only
through $W=\sum_js_j$, for all orders in a nonempty open set, if and
only if $A_1=\cdots=A_r$. For the \emph{unnormalized} profiles
$F_{1+A_j/q_j,s_j}(e^{q_jx})$, pure addition as a function of $\mu$
also requires $q_1=\cdots=q_r$.
\end{theorem}
\begin{proof}
Sufficiency in the normalized case is Theorem~\ref{clc:thm:addition}.
For necessity, weighted convolution is in $L^1$ as well as continuous,
so one can take its bilateral transform. If the function depends only
on $W$, its derivative in the direction
$\partial_{s_j}-\partial_{s_k}$ vanishes. By
\eqref{clc:eq:U-contour} and transform uniqueness, this gives
\[
 G_{\bq}(p)\prod_l(p+A_l)^{-s_l}
             [\log(p+A_k)-\log(p+A_j)]=0
\]
on the common vertical line. The prefactors never vanish there, so the
logarithms agree. Exponentiation implies $p+A_k=p+A_j$, hence
$A_k=A_j$. This holds for every pair.

In the unnormalized case there is an additional transform factor
$\prod_jq_j^{s_j}$. The same derivative now gives
$\log q_j-\log q_k-\log(p+A_j)+\log(p+A_k)=0$.
Thus $q_j(p+A_k)=q_k(p+A_j)$ on a line. Comparing the linear
coefficients in $p$ gives $q_j=q_k$, and then $A_j=A_k$.
\end{proof}

The statement concerns the entire convolution \emph{function} of $\mu$,
not a special numerical coincidence at one point or one order. It does
not rule out exceptional identities at $\mu=0$. Even with equal scales,
unequal centers already fail in the simplest example:
\begin{align*}
 \UU_{(1,1),(0,1)}((1,0);0)&=\zeta(2),\\
 \UU_{(1,1),(0,1)}((0,1);0)&=\zeta(2)-1.
\end{align*}
To verify these, the zero-order profile is independent of its center.
The two integrals therefore reduce respectively to the equal-center
weight-one formulas at $A=0$ and $A=1$.

\subsection{What commensurability really characterizes}
Consider the following explicitly restricted elementary class on $x<0$:
finite sums of functions
\begin{equation}\label{clc:eq:kernel-class}
 \frac{P_\nu(x)e^{\rho_\nu x}}
             {1-\epsilon_\nu e^{Qx}},
 \qquad Q>0,\quad |\epsilon_\nu|=1,
\end{equation}
where each $P_\nu$ is a polynomial, and the finite data $\rho_\nu$,
$\epsilon_\nu$ and $P_\nu$ are constant in $x$. The same $Q$ is used
for every term. Complex exponents are allowed.

\begin{theorem}[Common-step kernel rigidity]\label{clc:thm:commensurate}
For positive real scales $\bq$, the logistic convolution $K_{\bq}$
has a representation of the form \eqref{clc:eq:kernel-class} on $x<0$
if and only if all scale ratios are rational.
\end{theorem}
\begin{proof}
The sufficiency is \eqref{clc:eq:kernel}, since commensurate scales have
a positive common multiple $Q$.
For necessity, suppose such a finite representation exists and fix
$\delta>0$. On $x< -\delta$, expand the denominators as geometric
series. Integrating each term of its negative-half-line Laplace
transform gives meromorphic functions of $p$ whose poles lie in the
finite union
\[
 \bigcup_\nu\{\rho_\nu+nQ:n\ge0\}.
\]
Polynomial factors only increase pole multiplicities. The upper cutoff
$-\delta$ supplies exponential decay in $n$, so after omitting any
finite initial terms the integrated series is normally convergent on
compact subsets away from that pole set. This establishes the stated
meromorphic pole restriction, regardless of any cancellations between
terms of the finite representation.

The Laplace transform over $[-\delta,\infty)$ is holomorphic for
$\Rea p>0$. Indeed the weighted convolution estimates with arbitrarily
small positive weight show that $K_{\bq}(x)$ grows more slowly than
$e^{\varepsilon x}$ for every $\varepsilon>0$ as $x\to+\infty$.
The full transform agrees with $G_{\bq}(p)$ on $0<\Rea p<q_*$,
so meromorphic continuation forces every positive pole of $G_{\bq}$
to belong to the finite union above.

But $G_{\bq}$ has an uncancelled pole at every $nq_j$, $n\ge1$:
cosecant factors have no zeros, so another factor cannot remove it.
For a fixed $j$, two of the infinitely many $nq_j$ lie in the same
one of the finitely many progressions. Subtracting them gives
$(n-n')q_j=(k-k')Q$, hence $q_j/Q\in\mathbb Q$.
This holds for each $j$, proving commensurability.
\end{proof}

This theorem is a characterization of a fixed common-step elementary
kernel class. It is \emph{not} a claim that incommensurate integrals
cannot be expressed by other special functions, cannot have isolated
exact values, or yield arithmetically independent constants.
The next section supplies abundant isolated exact values for arbitrary
positive scales, emphasizing the distinction.
